\section{Representation Learning and Autoencoders}
dimensionality reduction, PCA, AE, latent space, VAE.

use papers on autoencoders and ML:
\begin{itemize}
    \item “Reducing the Dimensionality of Data with Neural Networks” — Geoffrey E. Hinton \& Ruslan R. Salakhutdinov (2006). General AE paper, comparison with PCA, deep AEs as a nonlinear dimensionality-reduction approach.
    \item “Representation Learning: A Review and New Perspectives” — Yoshua Bengio, Aaron Courville \& Pascal Vincent (2013). General representation learning review, motivation for learning useful features from high-dimensional data, comparison of representation-learning approaches.
    \item “Auto-Encoding Variational Bayes” — Diederik P. Kingma \& Max Welling (2013/2014). Foundational VAE paper, probabilistic latent representations, reparameterization trick and KL-divergence regularization.
    \item “Extracting and Composing Robust Features with Denoising Autoencoders” — Pascal Vincent, Hugo Larochelle, Yoshua Bengio \& Pierre-Antoine Manzagol (2008). Denoising AE paper, learning robust representations by reconstructing clean inputs from corrupted data, extension of the standard reconstruction objective. MAYBE, as the task used in thesis doesn't relate to denoising...
    \item “On Lines and Planes of Closest Fit to Systems of Points in Space” — Karl Pearson (1901). Foundational PCA work, linear dimensionality reduction, useful as the classical statistical comparison to autoencoder-based dimensionality reduction.
\end{itemize}
